\bookchapter{Chapitre 30 : Diffusion d'écran, présentations et connectivité automobile}

\chapterintro{L'écran du smartphone ne fait que 7 pouces au maximum, mais votre téléviseur en fait 55 et votre projecteur 100. « Projeter » l'image du smartphone sur un grand écran offre une expérience d'une tout autre dimension. Ce chapitre explique la duplication d'écran, la présentation et la connectivité automobile, trois scénarios fréquents, pour étendre les capacités du smartphone vers des écrans plus grands.}

\booksection{I. Les trois camps de la duplication d'écran sans fil — ne les confondez pas}

La duplication d'écran sans fil semble être toujours « l'image du smartphone sur la télé », mais il existe trois technologies sous-jacentes, incompatibles entre elles :

\begin{center}
\begin{tabularx}{\linewidth}{|X|X|X|}
\hline
\textbf{Technologie} & \textbf{Qui l'utilise} & \textbf{Caractéristiques} \tabularnewline
\hline
AirPlay & iPhone/iPad/Mac & Écosystème Apple, l'expérience la plus fluide. Nécessite Apple TV ou un téléviseur compatible AirPlay \tabularnewline
Miracast & Smartphones Android & Connexion directe Wi-Fi Direct, pas besoin de routeur. La « duplication sans fil » des smartphones Android l'utilise par défaut \tabularnewline
DLNA & Toutes plateformes & Pousse uniquement les fichiers multimédias (vidéo/musique/image), ne duplique pas l'écran. Idéal pour pousser une vidéo vers la télé \tabularnewline
\hline
\end{tabularx}
\end{center}

\booksubsection{Méthodes de fonctionnement par plateforme :}

\begin{itemize}
  \item \textbf{iPhone}: Centre de contrôle > Duplication d'écran > choisissez votre téléviseur/Apple TV. Si le téléviseur ne prend pas en charge AirPlay, vous pouvez acheter une Apple TV ou une clé TV compatible AirPlay.
  \item \textbf{Android}: Tirez le panneau de notifications > trouvez le bouton « duplication sans fil » / « duplication intelligente » / « interaction multi-écran » > choisissez le téléviseur. Le smartphone et le téléviseur doivent être sur le même Wi-Fi.
  \item \textbf{Solution universelle}: Si le téléviseur est trop ancien pour la duplication sans fil, branchez un Chromecast ou un boîtier Xiaomi/Huawei sur le port HDMI du téléviseur pour recevoir la duplication.
\end{itemize}

\booksection{II. Duplication filaire — plus stable, latence nulle}

La duplication sans fil est pratique mais a de la latence, ce qui gâche l'expérience en jeu et en présentation. La duplication filaire résout ce problème :

\begin{itemize}
  \item \textbf{iPhone}: Nécessite un câble adaptateur Lightning/USB-C vers HDMI (environ 100-200 yuans). Une fois branché, l'image du smartphone est directement envoyée au port HDMI du téléviseur/projecteur.
  \item \textbf{Android}: Certains smartphones prennent en charge la connexion directe USB-C vers HDMI (le smartphone doit prendre en charge DP Alt Mode, la plupart des phares le font). Achetez un câble USB-C vers HDMI.
  \item \textbf{Scénario salle de réunion}: Si vous faites souvent des présentations avec votre smartphone, gardez un câble adaptateur HDMI dans votre sac. Les téléviseurs/projecteurs de salle de réunion ont presque tous un port HDMI, il suffit de brancher, c'est bien plus fiable que la duplication sans fil.
\end{itemize}

\booksection{III. Le smartphone devient ordinateur — mode bureau}

Certains smartphones Android prennent en charge le « mode bureau » — connecté à un moniteur, le smartphone affiche une interface de bureau, utilisable avec souris et clavier :

\begin{itemize}
  \item \textbf{Samsung DeX}: La fonction phare des smartphones Samsung. Connecté à un moniteur, une interface de bureau apparaît, avec multi-fenêtres, l'expérience la plus proche d'un vrai ordinateur. Avec un clavier et une souris Bluetooth, vous pouvez voyager sans ordinateur portable.
  \item \textbf{Mode bureau Huawei}: Les séries Huawei Mate/Pura connectées à un moniteur offrent une fonction similaire. Prend en charge le multi-fenêtres, la gestion de fichiers et les applications de bureau.
  \item \textbf{Motorola Ready For}: Les phares Moto le prennent également en charge, similaire à DeX.
  \item \textbf{Scénario d'utilisation}: Bureautique légère (e-mails, documents), présentation PPT, vidéo sur grand écran. Pas adapté aux opérations Excel intensives ou au design professionnel.
\end{itemize}

\booksection{IV. Connectivité automobile — la bonne façon d'utiliser le smartphone en conduisant}

Manipuler le smartphone en conduisant est illégal, mais la navigation, la musique et les appels sont indispensables. Le système de connectivité automobile permet d'utiliser les fonctions principales du smartphone sur l'écran de la voiture, le smartphone restant de côté :

\begin{itemize}
  \item \textbf{Apple CarPlay}: Le premier choix des utilisateurs iPhone. Connecté par USB ou sans fil à l'écran de la voiture, une interface simplifiée de l'iPhone apparaît — navigation (Amap/Baidu Maps), musique, téléphone, messages, podcasts. Contrôle vocal Siri mains libres.
  \item \textbf{Android Auto}: La solution correspondante pour les utilisateurs Android. Fonctions similaires à CarPlay, prend en charge Google Maps, Spotify, etc. En Chine, certaines fonctions sont indisponibles en raison des restrictions des services Google.
  \item \textbf{Huawei HiCar}: Solution développée par Huawei, taux de support élevé parmi les véhicules nationaux. Outre la navigation et la musique, elle offre des fonctions spéciales comme les appels vidéo Huawei et la fusion d'images de stationnement.
  \item \textbf{Baidu CarLife}: Solution multiplateforme, utilisable sur iOS et Android. Moins de fonctions mais grande compatibilité, de nombreux vieux véhicules ne supportent que CarLife.
\end{itemize}

\begin{tcbwarn}{Rappel de sécurité}
Quelle que soit la solution de connectivité automobile, utilisez autant que possible le contrôle vocal en conduisant. Définissez la destination à l'avance, choisissez la playlist avant de rouler, ne manipulez pas l'écran de la voiture en roulant.

\end{tcbwarn}

\booksection{V. Duplication d'écran en réunion — plus de panique}

Les réunions sont le scénario le plus fréquent pour la duplication d'écran, et aussi le plus propice aux problèmes. Voici quelques conseils pratiques :

\begin{itemize}
  \item \textbf{Arrivez en avance pour tester}: N'attendez jamais le début de la réunion pour connecter la duplication. Arrivez 10 minutes en avance pour tester la connexion. Le taux d'échec de la duplication Wi-Fi est bien plus élevé que vous ne le pensez.
  \item \textbf{Préparez un plan B filaire}: Gardez un câble adaptateur HDMI dans votre sac. Quand la duplication sans fil ne fonctionne pas, le filaire est la bouée de sauvetage.
  \item \textbf{Désactiver l'aperçu des notifications}: Pendant la duplication, vos messages WeChat apparaîtront sur le grand écran de la réunion. Avant la duplication, désactivez « aperçu du contenu des notifications sur l'écran de verrouillage » dans les réglages, ou activez le « mode présentation » (certains smartphones le prennent en charge, masquant automatiquement les notifications pendant la duplication).
  \item \textbf{Utilisez « présentation de documents » plutôt que « manipulation du smartphone »}: Si vous présentez un PPT, enregistrez-le directement sur le smartphone et ouvrez-le en plein écran avec WPS, plutôt que de dupliquer tout le bureau du smartphone. C'est plus professionnel et plus fluide.
\end{itemize}

\begin{tcbpractice}{Exercice : effectuez une duplication du smartphone vers le téléviseur}
\begin{enumerate}
  \item Assurez-vous que le smartphone et la TV sont sur le même WiFi.
  \item Ouvrez une photo ou une vidéo sur le smartphone.
  \item Essayez la duplication sans fil (iPhone : duplication d'écran, Android : bouton de duplication sans fil).
  \item En cas d'échec, vérifiez si le téléviseur prend en charge le protocole correspondant, ou essayez le push DLNA (touchez le bouton « dupliquer » dans l'application Album/Vidéo).
  \item Une fois la duplication réussie, essayez d'autres opérations sur le smartphone pour voir si le grand écran affiche en synchronisation.
\end{enumerate}

\end{tcbpractice}

\begin{tcbsummary}{Résumé du chapitre}
\begin{itemize}
  \item Trois camps de duplication sans fil : AirPlay (Apple), Miracast (Android), DLNA (push multimédia multiplateforme)
  \item La duplication filaire est plus stable avec latence nulle, un câble HDMI est indispensable en réunion
  \item Samsung DeX / mode bureau Huawei transforme le smartphone en ordinateur de bureau léger
  \item Connectivité automobile : CarPlay / Android Auto / HiCar / CarLife, quatre solutions
  \item Règles d'or de la duplication en réunion : tester à l'avance, préparer un plan B, désactiver l'aperçu des notifications
\end{itemize}

\end{tcbsummary}
